\section{From surface assemblies to affine loop families}
\label{sec:surface-families}

\subsection{The actual input model}

Fix $m\ge1$ and $k\ge0$. Each listed patch is a compact connected surface
$F_v$ with nonempty boundary, described by orientability, genus or crosscap
number, and boundary count. A canonical free generating system has
$1-\chi(F_v)$ generators. Every generator acts on the common fibre
$\Z/m\Z$ by a translation
\[
 x\longmapsto x+\alpha_{v,i}(z),\qquad
 \alpha_{v,i}(z)=a_{v,i}+L_{v,i}z\pmod m.
\]
All expressions use the same $k$ parameters $z$. Local covers may be
disconnected; this is essential for a nontrivial connectivity optimization.

A seam joins two different available boundary ports, possibly in the same
patch. A port can be used at most once. Its base-circle degree is
$\epsilon_e\in\{1,-1\}$, and its fibre map is
$x\mapsto x+s_e(z)$, where $s_e$ is affine. The graph whose vertices are
patches and whose edges are seams is required to be connected. Additional
listed linear congruences on $z$ are permitted. The input does not
assert an embedding of these surfaces into a knot exterior.

For an orientable genus-$g$ surface with $b$ boundary components we use
generators $a_i,b_i$ and $c_1,\ldots,c_{b-1}$; the final peripheral word
is the inverse of the product of handle commutators and the $c_i$.
For a nonorientable surface, squares of the crosscap generators replace
the commutators. These words determine the peripheral translations
$p_{v,j}(z)$ by signed addition. No arbitrary proposed peripheral system
is trusted as a description of a surface.

\begin{lemma}[Affine seam equations]
\label{lem:seam-equations}
A seam from $(u,i)$ to $(v,j)$ with base degree $\epsilon_e$ and
translation $s_e$ is well-defined on the entire boundary preimage if and
only if
\[
 p_{u,i}(z)=\epsilon_e p_{v,j}(z)\pmod m.
\]
Consequently all seam conditions and all supplied linear restrictions
form one modular system $Az=b$.
\end{lemma}
\begin{proof}
Following a boundary loop before or after the seam gives respectively
the translations $s_e+p_{u,i}$ and $\epsilon_e p_{v,j}+s_e$.
Equality is precisely the displayed congruence. It is sufficient because
an equivariant bijection of the fibre extends to an isomorphism of the
corresponding boundary covers. Every expression is affine, so moving its
constant term to the right produces a linear congruence.
\end{proof}

\subsection{A spanning tree removes redundant seam variables from the loops}

Choose a rooted spanning tree of the patch graph. Set $t_o(z)=0$ at the
root and define $t_v$ along tree edges by adding or subtracting the
corresponding seam shifts. Thus $t_v$ is the sheet displacement from the
root fibre to the fibre in patch $v$. Store its affine coefficients reduced
modulo $m$, together with its parent edge and direction.

The loop list consists of every patch generator translation and, for
each non-tree seam $e:u\to v$, the affine displacement
\begin{equation}
 t_u(z)+s_e(z)-t_v(z).\label{eq:cycle-voltage}
\end{equation}
Conjugating a patch translation by the tree transport does not change it,
since translations commute. The list therefore has
\[
 r=\sum_v(1-\chi(F_v))+|E|-|V|+1
\]
entries. Relations imposed by the surface gluings need not make this a
free generating system; that is unnecessary for orbit computation.

\begin{theorem}[Component count for a compiled family]
\label{thm:surface-components}
For every feasible $z$, let $h_1(z),\ldots,h_r(z)$ be the compiled loop
translations. The assembled covering surface has exactly
\begin{equation}
 \kappa(z)=\gcd(m,h_1(z),\ldots,h_r(z))\label{eq:surface-count}
\end{equation}
connected components, with the empty-list convention $\gcd(m)=m$.
\end{theorem}
\begin{proof}
Every point in the assembled cover can be connected to a point of the
root fibre: first travel within its patch to the chosen fibre, then along
the tree. Two points in the root fibre are connected precisely when a
based loop in the assembled base sends one sheet to the other.
Patch loops and non-tree seam loops generate this action. Their image
is the additive subgroup of $\Z/m\Z$ generated by the $h_i$.
By Bezout's identity this subgroup consists exactly of the multiples of
the displayed gcd. Its orbits are its cosets, and there are precisely
that many. This proof also covers self-seams, the absence of non-tree
seams, disconnected local covers, and $m=1$.
\end{proof}

Combining this theorem with the arithmetic optimizer gives a deterministic
polynomial-bit-time procedure for the smallest possible component count
over the entire feasible assembly family. The returned assignment is
converted back to a concrete full-boundary assembly, using the input
schema of the earlier surface-gluing kernel~\cite{projectors}.
The certificate verifier recompiles the original surface data and checks
the dual arithmetic identities; it ignores the producer's stored
compilation. It therefore never accepts a modified loop list merely
because it accompanies an otherwise valid arithmetic witness.

\subsection{A nontrivial connectedness example}

Take an annulus with twelve sheets, trivial local monodromy, and glue its
two boundary preimages by a degree-$-1$ base map with fibre shift $z$.
Before the gluing there are twelve disjoint annuli. After it there are
$\gcd(12,z)$ components. The affine-family optimizer selects a unit shift
and certifies a connected torus. This is a genuine change in global
connectivity, not a reclassification of an already connected local block.

For a constrained example, give the annulus local monodromy $6$ and
impose $2z=8\pmod{12}$. The two feasible values are $4$ and $10$.
The resulting component count is
\[
 \gcd(12,6,z)=2
\]
for both. A single algebraic certificate proves that no connected member
exists. Literal enumeration is unnecessary even when twelve is replaced
by a binary integer with thousands of bits.

\subsection{Geometric and representation costs}

Let $N$ count the listed patches, generators, seam endpoints, and
constraint rows. Every affine expression that is present has $k+1$ listed entries.
We charge $k$, the listed dimensions, and all coefficient bit lengths as
in Section~\ref{sec:affine-cyclic-family}; $N$ need not count $N$
nonempty expressions, since a disc patch can have none. Constructing the
peripheral words and spanning tree, and adding their coefficient vectors,
takes $O(N(k+1))$ modular operations. The arithmetic solver charges its
own matrix dimensions and coefficient bit lengths. The parameter count
$k$ is an explicit dimension, not a constant-cost binary instruction to
construct a vector of astronomical length.

All loops remain translations. A negative base-boundary degree is not a
reflection of the sheet fibre and is not an orientation character for
the total surface. Allowing arbitrary signs or arbitrary affine units in
the fibre would change the algebra and is outside this compiled model.
The implementation preserves the base surface and seam directions so a
later classifier can evaluate the realized surface's genus and
orientability. Connectivity alone does not supply those answers.
